\documentclass{article}

% Packages
\usepackage[utf8]{inputenc}
\usepackage[T1]{fontenc}
\usepackage{times} % Times New Roman font
\usepackage[margin=3cm,left=3.5cm]{geometry} % Set margins
\usepackage{setspace} % For line spacing
\usepackage{graphicx} % For including figures
\usepackage[labelfont=bf]{caption} % Bold figure labels
%\usepackage{natbib} % For bibliography management
\usepackage{hyperref} % For hyperlinks
\usepackage{titling} % For custom title page
\usepackage{float} % For figure placement control
\usepackage[font=small,labelfont=bf]{caption}
\usepackage{tikz}
\usepackage[mode=buildnew]{standalone} 
\usepackage[style=apa, backend=biber]{biblatex}


\addbibresource{references.bib}

\usetikzlibrary{positioning, arrows.meta}

% set spelling to UK English
\usepackage[british]{babel}

% Line spacing (1.5)
\setstretch{1.5}

% Customize hyperref colors
\hypersetup{
    colorlinks=true,
    linkcolor=black,
    citecolor=black,
    urlcolor=blue
}

%\setlength{\parindent}{0pt}

\begin{document}

\begin{titlepage}
    \centering
    \vspace*{3cm}
    
    {\Huge\bfseries A Conservation-Driven Computational Screen for RNA Motifs\par}
    \vspace{2cm}
    
    {\Large Daniel Dachs\par}
    \vspace{1.5cm}
    
    {\large BIOC490\par}
    \vspace{1cm}
    
    {\large Supervisor: Paul Gardner\par}
    \vspace{1cm}
    
    {\large Department of Biochemistry\par}
    {\large University of Otago\par}
    \vspace{1cm}
    
    \vfill
\end{titlepage}


\section{Introduction}

RNAs play a wide range of roles in cellular processes, many of which have been harnessed for biotechnological applications, and are likely to have many more functions yet to be discovered.
This introduction provides background and rationale for the screening of novel functional RNA motifs in bacterial genomes, and outlines the aims and significance of the project.

\subsection{Ribonucleic Acids}

The central dogma of molecular biology, first articulated by Francis Crick, describes the flow of genetic information: DNA $\rightarrow$ RNA $\rightarrow$ protein \parencite{Crick1970}.
In this classical view, RNA serves primarily as an intermediate messenger, transcribing genetic information from DNA and translating it into proteins at the ribosome.
However, this perspective underrepresents the diversity of catalytic, regulatory, and structural roles played by RNA molecules.

In the hypothetical RNA world, RNA molecules are thought to have performed many of the roles now largely carried out by proteins and DNA, including catalysis, regulation, and information storage \parencite{Kruger1982, Guerrier-Takada1983, Gilbert1986, Breaker2012}.
Although many of these functions have since been supplanted by proteins and DNA, numerous RNA-based regulatory mechanisms persist in modern organisms.

\subsection{RNA Structure Determines Function}

RNA molecules, much like the polypeptides that make up proteins, fold into complex structures that enable diverse functions \parencite{Wang2021}.
They typically exist as single-stranded molecules, which can fold back on themselves through intramolecular base pairing to form secondary and tertiary structures.

RNA secondary structure arises from Watson--Crick base pairing (A--U and G--C) as well as wobble pairing between nucleotides within the same strand \parencite{Tinoco1999}. 
Common secondary structures include hairpins, consisting of a base-paired stem capped by an unpaired loop; internal loops and bulges, made up of unpaired nucleotides within a helical stem; junctions where three or more helical stems converge; and pseudoknots, where nucleotides in a loop pair with regions outside that loop \parencite{Staple2005}.

Beyond secondary structure, RNAs fold into three-dimensional structures through long-range interactions \parencite{Batey2004}.
These include base triplets, ribose zippers, A-minor motifs, and tetraloop--receptor interactions, which together stabilize compact RNA folds.

The relationship between RNA structure and function is exemplified by riboswitches, a class of regulatory RNAs.
The aptamer domain folds into a specific three-dimensional pocket that selectively binds a target ligand through a combination of electrostatic interactions and shape complementarity \parencite{Kavita2023}.
Ligand binding stabilizes one conformational state over another, thereby modulating transcription termination or initiation.
This allosteric mechanism—whereby binding at an aptamer site affects a distal functional region—demonstrates how RNA structure enables regulatory function.

\subsection{Functional RNAs in Biology and Biotechnology}

The diverse functional roles of RNA are driven by their structural versatility and chemical reactivity. 
Naturally occurring functional RNAs include autonomous RNA systems such as riboswitches and ribozymes, as well as RNAs that function in complexes with proteins.
In contrast, \emph{in vitro}-selected RNAs include fluorogens and aptazymes \parencite{Micura2020}.

\subsubsection*{Ribonucleoprotein Complexes}

While some RNAs function independently, many operate in ribonucleoprotein (RNP) complexes.
The ribosome is one such RNP with a critical function in all cells.
This large complex catalyses peptide bond formation, with ribosomal RNA (rRNA) forming the catalytic core, while ribosomal proteins provide structural support and regulatory functionality \parencite{Wahl2009}.
Spliceosomes are another example of essential RNPs in eukaryotes \parencite{Wahl2009}.

In bacteria, several classes of RNPs mediate post-transcriptional regulation.
Small regulatory RNAs (sRNAs) often require the RNA chaperone Hfq to facilitate base pairing with target mRNAs, thereby modulating transcript stability \parencite{Vogel2006}.

CRISPR--Cas systems exemplify the diversity and complexity of bacterial RNPs.
CRISPR RNAs (crRNAs) guide Cas nucleases to complementary DNA or RNA targets, enabling sequence-specific nucleic acid cleavage \parencite{Barrangou2007}.
This mechanism has famously been repurposed as a gene-editing tool \parencite{Jinek2012}.

\subsubsection*{Ribozymes}

Ribozymes are catalytic RNAs that accelerate chemical reactions, much like protein enzymes.
Beyond the ribosome and spliceosome, several small self-cleaving ribozymes exist in nature, including hammerhead, hairpin, hepatitis delta virus (HDV), glmS, twister, and pistol ribozymes \parencite{Roth2014, Weinberg2015, Micura2020}. 

\subsubsection*{Riboswitches}

Riboswitches are structured RNA domains that regulate gene expression in response to environmental stimuli, most commonly the presence of a metabolite \parencite{Breaker2023}.
Although seemingly less abundant than protein regulators that perform similar functions, riboswitches are now recognized as important regulatory elements across all domains of life.
They are typically embedded within messenger RNAs and control gene expression through ligand-induced conformational changes.

Riboswitch mechanisms vary widely.
Figure~\ref{fig:riboswitch-mechanism} illustrates several experimentally validated regulatory strategies.
Transcriptional riboswitches regulate the formation of Rho-independent terminators—intrinsic termination structures consisting of a hairpin followed by a uridine-rich tract.
Depending on the riboswitch, ligand binding may promote terminator formation or stabilize an antiterminator structure, enabling transcriptional read-through.
Similarly, translational riboswitches can sequester or expose the ribosome-binding site in response to ligand binding \parencite{Breaker2023}.

\subsubsection*{Biotechnological Applications}

RNA has been central to several major biotechnological breakthroughs.
CRISPR--Cas9 gene editing harnesses guide RNAs for programmable genome modification \parencite{Jinek2012}.
mRNA vaccines against COVID-19 demonstrated the therapeutic potential of synthetic mRNAs \parencite{Polack2020}.
RNA interference (RNAi) enables targeted gene silencing in both research and therapeutic contexts \parencite{Fire1998}.

RNA aptamers such as Spinach bind fluorogenic compounds and fluoresce upon binding, enabling RNA visualization in living cells \parencite{Paige2011}.
The photoriboswitch Werewolf-1, generated through \emph{in vitro} selection, responds to light by binding photoisomerizable ligands, offering potential for optogenetic control of gene expression \parencite{Rotstan2020}.
These RNAs, involved in fluorescence and photoreception, have broad relevance for medical research, biotechnology, and agricultural applications \parencite{Micura2020}.

\subsection{Current Approaches to Functional RNA Discovery}

It is predicted that a large number of riboswitches remain undiscovered, potentially on the order of thousands \parencite{Kavita2023}.
Ronald Breaker and colleagues have pioneered systematic riboswitch and ribozyme discovery, identifying over 55 distinct riboswitch classes to date \parencite{Breaker2022}.
Other major RNA-related discoveries, such as ribosomes and CRISPR--Cas systems, occurred through more exploratory research approaches \parencite{Brück2025}.

\subsubsection*{Experimental Approaches}

Transcriptome mapping approaches have proven successful for unbiased riboswitch discovery.
Term-seq maps RNA 3$'$ termini genome-wide to identify premature transcription termination sites characteristic of riboswitches and attenuators \parencite{Dar2016}.
This approach has uncovered numerous antibiotic-responsive riboswitches controlling resistance genes in bacterial pathogens and within the human microbiome.

RNA structure probing methods characterize secondary structures, tertiary interactions, and dynamic conformational changes.
SHAPE-seq (Selective 2$'$-Hydroxyl Acylation analysed by Primer Extension Sequencing) chemically modifies unpaired nucleotides, enabling genome-wide RNA structure mapping \parencite{Lucks2011}.
This approach has revealed structured regions across transcriptomes, including RNA switches, ribozymes, and protein-binding sites.
DMS-MaPseq extends this strategy to RNA structural probing \emph{in vivo} \parencite{Zubradt2017}.
A range of complementary methods also exist for studying protein--RNA interactions from both protein- and RNA-centric perspectives \parencite{Yi2025}.

While these experimental approaches are powerful for validation and mechanistic characterization, they are labour-intensive, may miss RNAs with unknown functions, and are difficult to scale across large numbers of organisms.

\subsubsection*{\emph{In Vitro} Selection and Computational Design}

\emph{In vitro} selection methods, such as SELEX, generate novel functional RNAs by applying selective pressure to large libraries of random sequences.
This technique was used in the generation of the photoriboswitch Werewolf-1 \parencite{Rotstan2020}.

Computational design approaches have also been used to generate synthetic functional RNAs, with tools such as RfamGen and RhoDesign \parencite{Sumi2024}.
However, many major RNA discoveries—including CRISPR--Cas systems, mRNA vaccine strategies, and numerous catalytic RNAs—were identified in nature rather than designed \emph{de novo} \parencite{Micura2020, Barrangou2007}.
Beyond biotechnology, the discovery of natural RNA-based regulatory systems provides important biological insights, as illustrated by the identification of fluoride-sensing riboswitches in bacteria \parencite{Breaker2012}.

\subsubsection*{Computational Approaches}

Bioinformatic sequence-based searches using tools such as BLAST can identify homologous sequences, but many functional RNAs exhibit limited sequence conservation and are detectable only through structural conservation.
Structure-based searches have proven particularly successful; for example, a bioinformatic screen by \textcite{Roth2014} identified approximately 2,700 twister ribozymes.
Such methods search for combined sequence and secondary structure patterns, but typically require existing alignments and structure models, limiting their utility for discovering entirely novel RNA classes.

RNA structure prediction remains a challenging problem, particularly when only a single sequence is available and no homologous structures are known.
Comparative approaches using multiple homologous sequences provide greater predictive power for RNA structure inference \parencite{Kagaya2026}.

Comparative genomics leverages evolutionary conservation of RNA structures across related species, enabling the identification of conserved structural features through covariation analysis.
These conserved features are referred to as RNA motifs—regions showing evidence of selective constraint on structure, indicative of functional importance.
Such motifs may correspond to riboswitches, ribozymes, or other classes of functional RNAs.

Zasha Weinberg's laboratory has had notable success using the tool CMfinder, which identifies conserved RNA secondary structures using a comparative genomics framework \parencite{Yao2006}.
More recently, \textcite{Sack2025} identified 16 novel structured RNA classes in plants using a CMfinder-based bioinformatic screen focused on alternative splicing regulation.

\textcite{Roth2014} also employed CMfinder to identify the twister ribozyme, now known to be widespread in nature.
In that study, hammerhead and twister ribozymes were frequently observed near one another and in association with particular protein-coding genes, though the biological significance of this pattern remains unclear.
\textcite{Weinberg2015} followed up on this observation by screening genomic regions near known ribozymes and ribozyme-associated protein domains, leading to the discovery of additional ribozyme classes, including pistol and hatchet.

\subsection{Gene Family-Targeted Searches}

Beyond genome-wide screens, gene family-targeted searches have proven effective for the discovery of riboswitches and other functional RNAs \parencite{Micura2020, Kavita2023}.
Many riboswitches are encoded upstream of specific gene families, reflecting shared regulatory mechanisms.
This targeted strategy is less computationally intensive than genome-wide approaches and enables broader analyses across more species.

Gene family-targeted searches exploit evolutionary conservation while incorporating the biological assumption that functionally related genes, particularly those sharing conserved protein domains, may also share conserved regulatory RNA elements.

Due to its relevance for biotechnology, this project focuses on screening conserved RNA structures located near light-sensing gene families.
Bacterial light-sensing systems often involve complex regulatory networks, and RNA-based mechanisms have been implicated in light acclimation responses in cyanobacteria \parencite{Hihara2015}.

A recent study by \textcite{Geissler2026} describes a bioinformatic screen similar in concept to the one undertaken here.
In that work, orthologous cyanobacterial genes involved in photosynthetic pathways were targeted for nearby conserved RNA structure discovery using CMfinder.
Approximately 400 novel conserved RNA structures were identified.

\subsection{Bacterial Genomes as a Search Space}

Riboswitches have been most extensively discovered and experimentally validated in bacteria and appear to be relatively rare in eukaryotes \parencite{Kavita2023}.
This makes bacterial genomes an optimal search space for riboswitch discovery efforts.

Bacterial genomes span extensive phylogenetic diversity.
As of April 2026, over 870,000 bacterial genomes representing approximately 180,000 species clusters have been sequenced and curated in databases such as the Genome Taxonomy Database (GTDB) \parencite{GTDB2026}.
Bacteria diverged more than three billion years ago, providing the deep evolutionary timescales required for detecting conserved RNA structures through covariation analysis.

Bacterial regulatory architectures are also generally simpler than those of eukaryotes, with fewer post-transcriptional processing steps and tighter coupling between transcription and translation.
These features simplify both computational analyses and downstream experimental validation.

\subsection{Aims}

This project aims to develop and apply a protein domain-associated comparative genomics pipeline to discover novel functional RNA motifs in bacterial genomes.

The specific aims are:

\begin{enumerate}
    \item Compile a list of candidate structural RNA motifs with potential regulatory or functional roles.
    \begin{itemize}
        \item Generate multiple sequence alignments of transcript regions, including 5$'$ untranslated regions, 5$'$ ends of coding sequences, and annotated intergenic regions, from defined gene families across phylogenetically diverse bacterial species
        \item Predict conserved RNA secondary structures using comparative genomics approaches that identify covariation patterns indicative of selective constraint
        \item Prioritise candidate motifs based on predicted structural complexity, strength of covariation, and phylogenetic distribution
    \end{itemize}

    \item Develop a reusable computational pipeline capable of systematically screening for conserved RNA structural motifs across diverse gene families and bacterial lineages.
\end{enumerate}

\subsection{Hypothesis}

Evolutionarily conserved structural motifs within transcripts associated with specific gene families represent functional regulatory elements maintained by natural selection.
These elements may employ diverse mechanisms, including ligand-dependent structural switching (riboswitches), catalytic self-cleavage (ribozymes), protein binding, or other structure-based regulatory strategies, to control gene expression.

% Bibliography
\newpage
\printbibliography

\end{document}